% !TEX program = pdflatex
% ============================================================================
% Paper A — Method, D2 (intervention response). DRAFT, protocol v2.1.
% Status: describes the protocol as implemented in scripts/d2_v2.py and
% scripts/analyze_d2.py (configs/d2v21_*.yaml). No results here.
%
% Compiles on its own (pdflatex). To move it into the paper: copy everything
% between \begin{document} and \end{document}, and move the macros below to
% the paper's preamble.
% ============================================================================
\documentclass[11pt,a4paper]{article}
\usepackage[utf8]{inputenc}
\usepackage[T1]{fontenc}
\usepackage{lmodern}
\usepackage[a4paper,margin=2.5cm]{geometry}
\usepackage{amsmath,amssymb}
\usepackage[hidelinks]{hyperref}

\providecommand{\din}{\Delta_{\mathrm{in}}}
\providecommand{\dout}{\Delta_{\mathrm{out}}}
\providecommand{\LR}{\mathrm{LR}}
\providecommand{\Gain}{\mathrm{Gain}}
\providecommand{\Dir}{\mathrm{Dir}}

\begin{document}

\section{Method}
\label{sec:method}

\subsection{Intervention response (D2)}
\label{sec:d2}

D1 asks whether a reliability signal \emph{ranks} the primitives associated with geometric failure
above the others. A high ranking score does not imply that acting on the selected primitives changes
the reconstruction in the way the signal suggests. D2 tests this directly: we intervene on the
primitives a signal selects, let the optimization respond, and compare the response with that of a
size-matched random selection.

\paragraph{Setup.}
Let $\mathcal{G}=\{g_i\}_{i=1}^{N}$ be a 3DGS model trained on views $\mathcal{V}_{\mathrm{tr}}$, and
let $\mathcal{V}_{\mathrm{ho}}$ be held-out views not used in training. A reliability signal assigns a
score $E_i$ to every primitive; by convention a higher score means \emph{less} reliable. For a
selection ratio $K$, the signal selects $R_E=\operatorname{top}_K(E)$, the $\lceil KN\rceil$ primitives with the
highest scores. Each signal is compared with $n_r$ random selections $R^{(j)}_{\mathrm{rand}}$,
$j=1,\dots,n_r$, of the same size, drawn on the same model.

\paragraph{Intervention and re-optimization.}
The intervention removes the selected primitives. The remaining model is then re-optimized from the
trained checkpoint with the same loss and with the optimizer state (Adam moments) of the
surviving primitives preserved. Densification and opacity reset are disabled and the position
learning rate is kept at its final value. As a result, no primitive is created or deleted during
re-optimization, and every surviving primitive can be compared one-to-one across runs. We record the
model at $T\in\mathcal{T}$ re-optimization steps. The snapshot at $T=0$ isolates the immediate effect
of the intervention. Later snapshots show how much of that effect the optimization compensates for.

\paragraph{References and noise floor.}
3DGS optimization is not deterministic: two re-optimizations of the same model with different view
orders already differ. We therefore run two \emph{null interventions}, i.e.\ re-optimizations without
removal. $\mathcal{G}^{\mathrm{null}}_a$ uses the same view order as every intervention and serves as
the paired reference. $\mathcal{G}^{\mathrm{null}}_b$ uses a different view order and gives the noise
floor. All changes below are measured against $\mathcal{G}^{\mathrm{null}}_a$ \emph{at the same step}
$T$, and the same quantity measured between $\mathcal{G}^{\mathrm{null}}_b$ and
$\mathcal{G}^{\mathrm{null}}_a$ is the noise floor. At $T=0$ both references equal the trained model
and the noise floor is zero.

\paragraph{Region of an intervention.}
Image-space locality requires a region in which a response is \emph{expected}. A geometric footprint
of the selected primitives (their projected extent) is too permissive: the footprint of a few percent of
the primitives covers most of the object, so every selection, random or not, looks local. We instead use
the share of each pixel's blending weight that the selection carries in the trained model,
\begin{equation}
  \omega_v(p) \;=\; \sum_{i\in R} w_{i,v}(p), \qquad
  w_{i,v}(p) = \alpha_{i,v}(p)\prod_{j\prec i}\bigl(1-\alpha_{j,v}(p)\bigr),
\end{equation}
where $\prec$ is front-to-back order along the ray through pixel $p$ of view $v\in\mathcal{V}_{\mathrm{ho}}$.
The map accounts for occlusion and opacity. It is obtained with one extra rendering pass that uses
the indicator of $R$ as the per-primitive feature. The region is
$M_R=\{(v,p):\omega_v(p)\ge\theta_w\}$. Selections differ widely in how much of the image they
dominate, so when $M_R$ covers less than a fraction $\rho$ of all held-out pixels we grow it with the
pixels of highest $\omega_v(p)$, up to that fraction. Without this floor, a selection whose region spans
a few pixels obtains an inflated signal-to-noise ratio simply because noise in a tiny region is small.

\paragraph{D2.1 — Locality.}
Let $\Delta(v,p)$ be the mean absolute RGB difference between the rendering of the intervened model and
that of $\mathcal{G}^{\mathrm{null}}_a$. With $\din$ and $\dout$ the mean of $\Delta$ inside and outside
$M_R$,
\begin{equation}
  \LR_{\mathrm{img}}(R) = \frac{\din}{\din+\dout}.
\end{equation}
Both terms are per-pixel means, so the ratio does not grow with region size. In 3D, let $\delta_i$ be
the displacement of the center of a surviving primitive between the intervened model and
$\mathcal{G}^{\mathrm{null}}_a$. Splitting survivors into those within distance $r$ of a removed
primitive and the rest gives $\LR_{\mathrm{3D}}$, defined in the same way. Locality alone is not
evidence that the selection is related to failure: a selection of influential but correct primitives can
also produce a local response. D2.2 addresses this.

\paragraph{D2.2 — Direction.}
A signal states that its top-ranked primitives are unreliable, which implies that removing them should
not degrade the reconstruction more than removing random primitives. Ideally, it should improve it. We
measure the change, relative to $\mathcal{G}^{\mathrm{null}}_a$, of
(i) held-out image error ($1-\mathrm{SSIM}$ and $L_1$ on $\mathcal{V}_{\mathrm{ho}}$), and
(ii) local geometric error with respect to the reference surface $\mathcal{S}$, restricted to the
neighborhood of radius $r$ around the removed primitives. Local geometric error is split into
\begin{equation}
  \mathrm{Acc} = \operatorname*{mean}_{\mu\in\mathcal{M}_{\mathrm{loc}}} d(\mu,\mathcal{S}), \qquad
  \mathrm{Comp} = \operatorname*{mean}_{s\in\mathcal{S}_{\mathrm{loc}}} d(s,\mathcal{M}),
\end{equation}
where $\mathcal{M}$ are the centers of primitives with opacity above $0.1$ and the subscript
$\mathrm{loc}$ restricts to the neighborhood. Accuracy degrades when erroneous primitives
(e.g.\ floaters) remain. Completeness degrades when valid surface is removed and leaves holes. The two
can move in opposite directions under the same intervention, and only their split shows whether a
selection removed erroneous primitives or valid surface. For each quantity $q$,
$\Dir_q=-(q_{\mathrm{int}}-q_{\mathrm{null}_a})$, so that positive values mean the intervention improved
the reconstruction.

\paragraph{D2.3 — Matched contrast.}
Every quantity is reported as a paired difference with the random selections of the same model,
\begin{equation}
  \Gain_X(E) = X(R_E) - \frac{1}{n_r}\sum_{j=1}^{n_r} X\bigl(R^{(j)}_{\mathrm{rand}}\bigr),
  \qquad X\in\{\LR_{\mathrm{img}},\LR_{\mathrm{3D}},\Dir_{1-\mathrm{SSIM}},\Dir_{\mathrm{Acc}},\Dir_{\mathrm{Comp}}\},
\end{equation}
together with the response-to-noise ratio
$\mathrm{SNR}=\din(\text{intervention vs.\ null}_a)\,/\,\din(\text{null}_b\text{ vs.\ null}_a)$, computed in
the same region $M_{R_E}$. Positive gains mean the signal's selection is more local, or improves the
reconstruction more, than random selection.

\paragraph{Outcomes.}
A signal is \emph{measurable} if its SNR is at least $2$ in every seed at the primary step. Gains of
signals that are not measurable are not interpreted. Removing primitives that do not contribute to
the image trivially ``beats'' random removal, which does change the image. For a measurable signal,
each gain is classified as \emph{expected} if it has the same positive sign in every seed and its mean
is at least two standard errors from zero, \emph{opposite} under the same condition with a negative
sign, and \emph{inconsistent} otherwise. This yields three outcomes per signal:
(i) the selection produces an expected locality gain and at least one expected direction gain, with no
opposite direction gain, which is evidence consistent with the signal's stated meaning;
(ii) the signal is not measurable, so its selection has little influence on the reconstruction;
(iii) at least one direction gain is opposite, so the stated meaning is not supported.
With the few seeds of the pilot experiments, this rule is a pre-specified heuristic, not a significance
test. The full study replaces it with paired Wilcoxon signed-rank tests against random selections,
with Holm--Bonferroni correction across signals (Sec.~\ref{sec:stats}).

These outcomes describe the response to the tested intervention under the stated settings. They do not
amount to a causal identification of which primitives cause the failure.

\paragraph{Implementation details.}
We use $K=5\%$, removal as the intervention, $\mathcal{T}=\{0,200,1000\}$ with $T=1000$ as the primary
step, $n_r=5$ random selections, $\theta_w=0.2$, $\rho=1\%$ of held-out pixels and $r=0.05$ scene
units. Signals are computed once on the trained model. The rasterizer is a pure-PyTorch
implementation of the 3DGS forward model (EWA projection with a $0.3$\,px low-pass filter, $3\sigma$
footprints, front-to-back compositing with $\alpha\le0.99$ and contributions below $1/255$ discarded).
It enumerates only primitive--pixel pairs with non-negligible opacity, and it matches a dense reference
implementation to within $10^{-6}$ relative error in both images and gradients. Unlike the CUDA
rasterizer, it does not terminate rays early and it sorts per pixel rather than per tile.
% TODO(Huy): runtime per re-optimization (≈3.4 min / 1000 steps, 160 px, ~28k primitives, GTX 1650);
%            move the rasterizer paragraph to a shared "Implementation" section if D1/D3 use it too.
% TODO(Huy): opacity scaling and position perturbation as additional interventions (Phase 1).

\subsection{Statistical analysis}
\label{sec:stats}
% TODO(Huy): unit of analysis (scene, view configuration, seed); pre-specified primary metrics;
%            paired Wilcoxon signed-rank vs. random with Holm--Bonferroni; bootstrap 95% CIs.
\emph{To be written.}

\end{document}
